\documentclass[12pt,a4paper]{article}
\usepackage[utf8]{inputenc}
\usepackage[T1]{fontenc}
\usepackage[brazil]{babel}
\usepackage{helvet}
\renewcommand{\familydefault}{\sfdefault}
\usepackage{geometry,setspace,longtable,array,indentfirst,hyperref,enumitem,ragged2e}
\usepackage[table]{xcolor}
\usepackage{titlesec,fancyhdr,xurl}

\definecolor{AzulCorporativo}{RGB}{0,82,155}
\definecolor{AzulPetroleo}{RGB}{23,104,133}
\definecolor{CinzaEscuro}{RGB}{64,64,64}
\definecolor{LaranjaDestaque}{RGB}{230,115,0}
\hypersetup{
	colorlinks=true,
	linkcolor=AzulPetroleo,
	urlcolor=AzulCorporativo,
	citecolor=LaranjaDestaque,
	pdftitle={Modelagem Física - Modelo Analítico de Dados para E-commerce},
	pdfauthor={Lucas Dias Noronha}
}
\titleformat{\section}{\color{AzulCorporativo}\normalfont\Large\bfseries}{\thesection}{1em}{}
\titleformat{\subsection}{\color{AzulPetroleo}\normalfont\large\bfseries}{\thesubsection}{1em}{}
\titleformat{\subsubsection}{\color{CinzaEscuro}\normalfont\normalsize\bfseries}{\thesubsubsection}{1em}{}
\geometry{left=3cm,right=2cm,top=3cm,bottom=2cm,headheight=45pt}
\onehalfspacing
\setlength{\parskip}{0.35em}
\setlength{\emergencystretch}{2em}
\pagestyle{fancy}
\fancyhf{}
\renewcommand{\headrulewidth}{0pt}
\fancyhead[C]{\colorbox{AzulCorporativo}{\makebox[\dimexpr\textwidth-2\fboxsep\relax][l]{\textcolor{white}{\textbf{\ \ Modelagem Física}}}}}
\fancyfoot[C]{\color{CinzaEscuro}\thepage}
\newcommand{\projeto}{Modelo Analítico de Dados para E-commerce}
\newcommand{\autor}{Lucas Dias Noronha}
\newcommand{\dataDocumento}{\today}

\begin{document}
	\thispagestyle{empty}
	\vspace*{\fill}
	\begin{center}
		{\Huge \textcolor{AzulCorporativo}{\textbf{Modelagem Física}}}\\[0.5cm]
		{\Large \textcolor{AzulPetroleo}{\projeto}}\\[0.8cm]
	\end{center}
	\vspace*{\fill}
	\noindent
	\textbf{\textcolor{AzulCorporativo}{Autor:}} \autor
	\hfill
	\textbf{\textcolor{AzulCorporativo}{Data:}} \dataDocumento
	\newpage
	\setcounter{tocdepth}{1}
	\tableofcontents
	\newpage

	\section{Introdução}
	Este documento consolida a modelagem física do projeto \textit{\projeto}.
	O modelo lógico aprovado é materializado no PostgreSQL por meio de tipos
	concretos, nulabilidade, chaves, restrições, índices e scripts SQL
	reproduzíveis.

	A solução segue uma arquitetura ELT em três camadas:
	\begin{center}
		\textbf{Dataset Olist $\rightarrow$ RAW $\rightarrow$ CORE
		$\rightarrow$ ANALYTICS $\rightarrow$ Power BI / SQL}
	\end{center}

	A modelagem preserva a semântica e a granularidade aprovadas. Decisões de
	carga, transformação e consumo analítico são descritas por sua finalidade,
	sem depender da numeração operacional utilizada no gerenciamento do projeto.

	\section{Sistema Gerenciador de Banco de Dados}
	O SGBD adotado é o \textbf{PostgreSQL 18}, sempre em sua revisão secundária
	mais recente disponível. A escolha considera:
	\begin{itemize}[leftmargin=1.2cm]
		\item suporte nativo a chaves, integridade referencial, transações e
		restrições de domínio;
		\item tipos adequados para valores monetários exatos, coordenadas e dados
		temporais;
		\item recursos SQL para junções, agregações, funções de janela e CTEs;
		\item ferramentas de linha de comando para automação e validação;
		\item integração com Python, DBeaver e ferramentas de BI;
		\item manutenção ativa, documentação oficial e licença livre.
	\end{itemize}

	O DBeaver é utilizado como cliente gráfico para conexão, execução de scripts e
	inspeção de objetos. Ele não substitui o servidor PostgreSQL nem os testes
	reproduzíveis. Credenciais e configurações locais não são versionadas.

	\section{Arquitetura Física em Camadas}
	\subsection{RAW}
	O schema \texttt{raw} recebe uma representação fiel dos nove arquivos CSV.
	Há uma tabela por arquivo, os nomes originais das colunas são preservados e os
	valores de origem são armazenados como \texttt{text}. Dessa forma, uma falha
	de conversão não impede a ingestão nem elimina a evidência original.

	Cada tabela acrescenta somente:
	\begin{itemize}[leftmargin=1.2cm]
		\item \texttt{\_id\_raw}: identidade técnica da ocorrência;
		\item \texttt{\_arquivo\_origem}: nome do arquivo recebido;
		\item \texttt{\_carregado\_em}: instante da ingestão.
	\end{itemize}
	Os campos da fonte permanecem anuláveis e não recebem regras de negócio. A
	chave técnica não afirma unicidade sobre o conteúdo original.

	\subsection{CORE}
	O schema \texttt{core} materializa as nove tabelas relacionais aprovadas.
	Renomeações, conversões, deduplicação e validações ocorrem no fluxo
	RAW--CORE. O CORE é a fonte estruturada e íntegra para consultas e para a
	produção das estruturas analíticas.

	\subsection{ANALYTICS}
	O schema \texttt{analytics} é inicialmente criado sem objetos. Views, marts e
	métricas serão derivados do CORE quando suas granularidades, regras e
	finalidades de consumo estiverem definidas. Power BI e consumidores SQL
	deverão utilizar preferencialmente essa camada.

	\subsection{Mapeamento dos Arquivos para a RAW}
	{\scriptsize
	\renewcommand{\arraystretch}{1.25}
	\begin{longtable}{|>{\raggedright\arraybackslash}p{7.2cm}|>{\raggedright\arraybackslash}p{7.2cm}|}
		\hline
		\rowcolor{AzulCorporativo}
		\textcolor{white}{\textbf{Arquivo}} & \textcolor{white}{\textbf{Tabela RAW}} \\ \hline
		\endfirsthead
		\hline
		\rowcolor{AzulCorporativo}
		\textcolor{white}{\textbf{Arquivo}} & \textcolor{white}{\textbf{Tabela RAW}} \\ \hline
		\endhead
		\texttt{olist\_customers\_dataset.csv} & \texttt{raw.olist\_customers} \\ \hline
		\texttt{olist\_geolocation\_dataset.csv} & \texttt{raw.olist\_geolocation} \\ \hline
		\texttt{olist\_order\_items\_dataset.csv} & \texttt{raw.olist\_order\_items} \\ \hline
		\texttt{olist\_order\_payments\_dataset.csv} & \texttt{raw.olist\_order\_payments} \\ \hline
		\texttt{olist\_order\_reviews\_dataset.csv} & \texttt{raw.olist\_order\_reviews} \\ \hline
		\texttt{olist\_orders\_dataset.csv} & \texttt{raw.olist\_orders} \\ \hline
		\texttt{olist\_products\_dataset.csv} & \texttt{raw.olist\_products} \\ \hline
		\texttt{olist\_sellers\_dataset.csv} & \texttt{raw.olist\_sellers} \\ \hline
		\texttt{product\_category\_name\_translation.csv} & \texttt{raw.product\_category\_name\_translation} \\ \hline
	\end{longtable}}

	\section{Evidências para os Tipos Físicos}
	Os CSVs foram perfilados como texto para evitar inferência indevida. Foram
	avaliados nulidade, comprimento, domínio, mínimo, máximo, escala decimal,
	formato temporal, unicidade e integridade referencial.

	{\small
	\renewcommand{\arraystretch}{1.25}
	\begin{longtable}{|>{\raggedright\arraybackslash}p{3.3cm}|>{\raggedright\arraybackslash}p{5.4cm}|>{\raggedright\arraybackslash}p{5.5cm}|}
		\hline
		\rowcolor{AzulCorporativo}
		\textcolor{white}{\textbf{Domínio}} & \textcolor{white}{\textbf{Evidência}} & \textcolor{white}{\textbf{Decisão}} \\ \hline
		\endfirsthead
		\hline
		\rowcolor{AzulCorporativo}
		\textcolor{white}{\textbf{Domínio}} & \textcolor{white}{\textbf{Evidência}} & \textcolor{white}{\textbf{Decisão}} \\ \hline
		\endhead
		Identificadores & 32 caracteres hexadecimais & \texttt{varchar(32)}; não converter para UUID \\ \hline
		Prefixo de CEP & exatamente cinco dígitos & \texttt{char(5)} com validação \\ \hline
		UF & dois caracteres; conjunto das UFs brasileiras & \texttt{char(2)} com domínio \\ \hline
		Cidades & maior ocorrência com 40 caracteres & \texttt{varchar(50)} \\ \hline
		Categoria & maior ocorrência com 46 caracteres & \texttt{varchar(50)} \\ \hline
		Valores monetários & duas casas; máximo 13.664,08 & \texttt{numeric(12,2)} \\ \hline
		Datas e horas & formato válido, sem deslocamento UTC & \texttt{timestamp without time zone} \\ \hline
		Coordenadas & precisão variável e pontos fora do Brasil & \texttt{double precision}; limites globais \\ \hline
	\end{longtable}}

	\section{Dicionário Físico do CORE}
	As colunas sem indicação de nulabilidade são obrigatórias. Nenhum valor de
	negócio recebe padrão implícito. Somente
	\texttt{geolocalizacao.id\_geolocalizacao}, inexistente na fonte, é gerado
	como \texttt{bigint generated always as identity}.

	{\small
	\renewcommand{\arraystretch}{1.18}
	\begin{longtable}{|>{\raggedright\arraybackslash}p{2.7cm}|>{\raggedright\arraybackslash}p{3.5cm}|>{\raggedright\arraybackslash}p{3.4cm}|>{\raggedright\arraybackslash}p{4.6cm}|}
		\hline
		\rowcolor{AzulCorporativo}
		\textcolor{white}{\textbf{Tabela}} & \textcolor{white}{\textbf{Coluna}} & \textcolor{white}{\textbf{Tipo}} & \textcolor{white}{\textbf{Integridade e observações}} \\ \hline
		\endfirsthead
		\hline
		\rowcolor{AzulCorporativo}
		\textcolor{white}{\textbf{Tabela}} & \textcolor{white}{\textbf{Coluna}} & \textcolor{white}{\textbf{Tipo}} & \textcolor{white}{\textbf{Integridade e observações}} \\ \hline
		\endhead
		cliente & id\_cliente & varchar(32) & PK; hexadecimal \\ \hline
		cliente & id\_cliente\_unico & varchar(32) & hexadecimal \\ \hline
		cliente & prefixo\_cep & char(5) & FK; cinco dígitos \\ \hline
		cliente & cidade / estado & varchar(50) / char(2) & texto não vazio; UF válida \\ \hline
		pedido & id\_pedido & varchar(32) & PK; hexadecimal \\ \hline
		pedido & id\_cliente & varchar(32) & FK e UNIQUE \\ \hline
		pedido & status\_pedido & varchar(20) & domínio dos estados observados \\ \hline
		pedido & datas do ciclo & timestamp without time zone & aprovação, envio e entrega podem ser nulos \\ \hline
		item\_pedido & id\_pedido, id\_item & varchar(32), smallint & PK composta; item positivo \\ \hline
		item\_pedido & id\_produto, id\_vendedor & varchar(32) & FKs \\ \hline
		item\_pedido & preço e frete & numeric(12,2) & valores não negativos \\ \hline
		produto & id\_produto & varchar(32) & PK; hexadecimal \\ \hline
		produto & nome\_categoria & varchar(50) & anulável \\ \hline
		produto & métricas e dimensões & smallint / integer & anuláveis e não negativas \\ \hline
		vendedor & id\_vendedor & varchar(32) & PK; hexadecimal \\ \hline
		vendedor & prefixo, cidade, estado & char(5), varchar(50), char(2) & FK e domínios geográficos \\ \hline
		pagamento & id\_pedido, sequencial & varchar(32), smallint & PK composta; FK; sequência positiva \\ \hline
		pagamento & tipo / parcelas / valor & varchar(20), smallint, numeric(12,2) & domínios e valores não negativos \\ \hline
		avaliacao & id\_avaliacao, id\_pedido & varchar(32) & PK composta; FK \\ \hline
		avaliacao & nota / datas & smallint, timestamp & nota entre 1 e 5 \\ \hline
		prefixo\_cep & prefixo\_cep & char(5) & PK; cinco dígitos \\ \hline
		geolocalizacao & id\_geolocalizacao & bigint identity & PK gerada \\ \hline
		geolocalizacao & prefixo\_cep & char(5) & FK \\ \hline
		geolocalizacao & latitude / longitude & double precision & limites globais \\ \hline
		geolocalizacao & cidade / estado & varchar(50) / char(2) & texto não vazio; UF válida \\ \hline
	\end{longtable}}

	As exclusões referenciadas usam \texttt{ON DELETE RESTRICT}; atualizações de
	identificadores usam \texttt{ON UPDATE NO ACTION}. Não são impostas
	comparações universais entre datas. Nulos descritivos e físicos observados em
	produtos são preservados, sem preenchimento artificial.

	\section{Estratégia Inicial de Índices}
	O PostgreSQL cria índices B-tree para PKs e restrições \texttt{UNIQUE}, mas
	não indexa automaticamente as colunas que referenciam FKs. Uma FK é considerada
	coberta quando suas colunas formam o prefixo à esquerda de um índice existente.

	Foram aprovados seis índices adicionais:
	\begin{itemize}[leftmargin=1.2cm]
		\item \texttt{idx\_cliente\_prefixo\_cep};
		\item \texttt{idx\_vendedor\_prefixo\_cep};
		\item \texttt{idx\_item\_pedido\_id\_produto};
		\item \texttt{idx\_item\_pedido\_id\_vendedor};
		\item \texttt{idx\_avaliacao\_id\_pedido};
		\item \texttt{idx\_geolocalizacao\_prefixo\_cep}.
	\end{itemize}

	Não há índices adicionais na RAW antes da implementação das transformações,
	pois eles aumentariam o custo de ingestão sem padrão de acesso comprovado.
	Também não há índices em ANALYTICS enquanto o schema estiver vazio. Depois da
	carga, estatísticas, consultas reais e planos de execução deverão orientar
	qualquer revisão.

	\section{Scripts e Operação}
	Os artefatos técnicos são:
	\begin{itemize}[leftmargin=1.2cm]
		\item \texttt{create\_schema.sql}: cria os três schemas e as tabelas;
		\item \texttt{create\_indexes.sql}: cria os seis índices adicionais;
		\item \texttt{validate\_schema.sql}: valida catálogo e comportamento;
		\item \texttt{drop\_indexes.sql}: remove os índices adicionais;
		\item \texttt{drop\_schema.sql}: remove tabelas e schemas em ordem segura.
	\end{itemize}

	No DBeaver, a criação é realizada executando
	\path{create_schema.sql}, \path{create_indexes.sql} e
	\path{validate_schema.sql}, nessa ordem. A criação falha
	intencionalmente quando os schemas já existem, protegendo objetos contra
	redefinição silenciosa. A remoção deve ser executada somente em ambiente
	descartável ou mediante autorização explícita.

	\section{Validação da Arquitetura}
	A arquitetura foi executada em uma instância limpa do PostgreSQL 18.3.
	O resultado consolidado foi:

	{\small
	\renewcommand{\arraystretch}{1.2}
	\begin{longtable}{|>{\raggedright\arraybackslash}p{8.5cm}|>{\centering\arraybackslash}p{2.5cm}|>{\centering\arraybackslash}p{2.5cm}|}
		\hline
		\rowcolor{AzulCorporativo}
		\textcolor{white}{\textbf{Verificação}} & \textcolor{white}{\textbf{Esperado}} & \textcolor{white}{\textbf{Obtido}} \\ \hline
		\endfirsthead
		\hline
		\rowcolor{AzulCorporativo}
		\textcolor{white}{\textbf{Verificação}} & \textcolor{white}{\textbf{Esperado}} & \textcolor{white}{\textbf{Obtido}} \\ \hline
		\endhead
		Schemas físicos & 3 & 3 \\ \hline
		Tabelas RAW / CORE / ANALYTICS & 9 / 9 / 0 & 9 / 9 / 0 \\ \hline
		Colunas RAW / CORE & 79 / 50 & 79 / 50 \\ \hline
		Colunas identity & 10 & 10 \\ \hline
		Chaves primárias / estrangeiras / UNIQUE & 18 / 9 / 1 & 18 / 9 / 1 \\ \hline
		Índices RAW / CORE / adicionais CORE & 9 / 16 / 6 & 9 / 16 / 6 \\ \hline
	\end{longtable}}

	A RAW aceitou deliberadamente valores fora dos domínios de negócio. O CORE
	rejeitou CEP inválido, chave duplicada e referência órfã. Os registros de teste
	foram revertidos. Uma segunda criação foi rejeitada sem danificar as tabelas
	existentes, e a remoção pôde ser repetida sem erro.

	\section{Conclusão}
	A modelagem física está consistente com os modelos conceitual e lógico e
	pronta para receber o processo ELT. A etapa seguinte deverá extrair os CSVs,
	carregá-los com rastreabilidade na RAW, transformar e validar os dados no CORE
	e reconciliar volumes e rejeições. Em seguida, estruturas analíticas derivadas
	do CORE poderão ser disponibilizadas para Power BI e consultas SQL.

	\section{Referências}
	\begin{itemize}[leftmargin=1.2cm]
		\item PostgreSQL. \textit{Versioning Policy}. Disponível em:
		\url{https://www.postgresql.org/support/versioning/}.
		\item PostgreSQL. \textit{PostgreSQL 18 Documentation}. Disponível em:
		\url{https://www.postgresql.org/docs/18/}.
		\item DBeaver. \textit{PostgreSQL driver documentation}. Disponível em:
		\url{https://dbeaver.com/docs/dbeaver/Database-driver-PostgreSQL/}.
	\end{itemize}
\end{document}
